% LaTeX-Vorlage für Spanisch-Dialoge (Südamerikanisches Spanisch)
% Diese Vorlage kann als Basis für alle Dialoge verwendet werden
% WICHTIG: Verwendung von südamerikanischem Spanisch - kein "vosotros"

\documentclass[12pt,a4paper]{article}
\usepackage[utf8]{inputenc}
\usepackage[spanish,german]{babel}
\usepackage[T1]{fontenc}
\usepackage{geometry}
\geometry{a4paper, margin=2.5cm}
\usepackage{titlesec}
\usepackage{enumitem}
\usepackage{multicol}
\usepackage{csquotes}

% Titel-Formatierung
\titleformat{\section}
{\Large\bfseries}
{}
{0em}
{}[\titlerule]

\titleformat{\subsection}
{\large\bfseries}
{}
{0em}
{}

% Abstände anpassen
\setlength{\parindent}{0pt}
\setlength{\parskip}{0.5em}

\begin{document}

% ============================================
% METADATEN
% ============================================
\title{\textbf{DIALOG: [Titel des Dialogs]}}
\author{Spanisch A1/A2}
\date{}
\maketitle

\textbf{Tema:} [Thema] \\
\textbf{Nivel:} A1/A2 \\
\textbf{Complejidad:} [Einfach/Leicht-Mittel/Mittel] \\
\textbf{Palabras:} [Anzahl der Wörter]

\vspace{1em}

\textit{Nota: Se recomienda revisar primero el vocabulario antes de leer el diálogo.}

\vspace{1em}

% ============================================
% DIALOG
% ============================================
\section*{Diálogo}

[HIER KOMMT DER DIALOG - 500-1000 Wörter]

% Beispiel-Struktur für Dialog:
% \textbf{Persona 1:} Hola, ¿cómo estás?
% 
% \textbf{Persona 2:} Muy bien, gracias. ¿Y tú?
% 
% \textbf{Persona 1:} También estoy bien.

\newpage

% ============================================
% VOCABULARIO
% ============================================
\section*{Vocabulario}

\begin{multicols}{2}
\begin{itemize}[leftmargin=*]
    \item \textbf{palabra} - Übersetzung
    \item \textbf{frase} - Übersetzung der Phrase
    % Weitere Vokabeln hier einfügen
\end{itemize}
\end{multicols}

\newpage

% ============================================
% PREGUNTAS DE COMPRENSIÓN
% ============================================
\section*{Preguntas de Comprensión}

\begin{enumerate}[leftmargin=*]
    \item [Frage 1]
    \item [Frage 2]
    % ... bis 20 Fragen
\end{enumerate}

\newpage

% ============================================
% VERBO IRREGULAR
% ============================================
\section*{Verbo Irregular: [VERB] ([deutsche Bedeutung])}

\textit{Nota: \enquote{vosotros/as} se usa solo en España, no en español sudamericano. Se incluye aquí para completitud.}

\begin{multicols}{2}

\subsection*{Presente}
\begin{itemize}[leftmargin=*]
    \item yo - [konjugiert]
    \item tú - [konjugiert]
    \item él/ella/usted - [konjugiert]
    \item nosotros/as - [konjugiert]
    \item vosotros/as - [konjugiert]
    \item ustedes - [konjugiert]
    \item ellos/ellas - [konjugiert]
\end{itemize}

\subsection*{Pretérito Indefinido}
\begin{itemize}[leftmargin=*]
    \item yo - [konjugiert]
    \item tú - [konjugiert]
    \item él/ella/usted - [konjugiert]
    \item nosotros/as - [konjugiert]
    \item vosotros/as - [konjugiert]
    \item ustedes - [konjugiert]
    \item ellos/ellas - [konjugiert]
\end{itemize}

\subsection*{Pretérito Imperfecto}
\begin{itemize}[leftmargin=*]
    \item yo - [konjugiert]
    \item tú - [konjugiert]
    \item él/ella/usted - [konjugiert]
    \item nosotros/as - [konjugiert]
    \item vosotros/as - [konjugiert]
    \item ustedes - [konjugiert]
    \item ellos/ellas - [konjugiert]
\end{itemize}

\columnbreak

\subsection*{Futuro Simple}
\begin{itemize}[leftmargin=*]
    \item yo - [konjugiert]
    \item tú - [konjugiert]
    \item él/ella/usted - [konjugiert]
    \item nosotros/as - [konjugiert]
    \item vosotros/as - [konjugiert]
    \item ustedes - [konjugiert]
    \item ellos/ellas - [konjugiert]
\end{itemize}

\subsection*{Pretérito Perfecto}
\begin{itemize}[leftmargin=*]
    \item yo - [konjugiert]
    \item tú - [konjugiert]
    \item él/ella/usted - [konjugiert]
    \item nosotros/as - [konjugiert]
    \item vosotros/as - [konjugiert]
    \item ustedes - [konjugiert]
    \item ellos/ellas - [konjugiert]
\end{itemize}

\end{multicols}

\newpage

% ============================================
% EXPLICACIÓN GRAMATICAL
% ============================================
\section*{Explicación Gramatical}

\textbf{[Titel der grammatikalischen Erklärung]}

[Hier kommt die Erklärung auf Deutsch]

\textbf{Ejemplos del texto:}
\begin{itemize}[leftmargin=*]
    \item "[Beispiel 1 aus dem Dialog]"
    \item "[Beispiel 2 aus dem Dialog]"
\end{itemize}

\end{document}
